\section{The Symbolism of the Horse}

\begin{quotex}
Try as I might, I never made my wife happy. So I made it my task to learn how to please a woman. I then was able to please a dozen women, although it never made me happy. So I became an itinerant horse trader.
\flright{\textsc{Mullah Nasruddin}}
\end{quotex}

Guenon tells us in \emph{Symbols of Sacred Science} that the Horse represents the Sun or Agni, and hence is a solar symbol. He also makes reference to this passage from Revelation 19:14-15:

\begin{quotex}
And the armies that are in heaven followed him Christ on white horses, clothed in fine linen, white and clean. And out of his mouth goes a sharp two-edged sword, that with it, he may strike the nations.
\end{quotex}


Similarly, according to the Puranas, Kalki, the final avatar, will also come riding on a white horse at the end of the Kali Yuga, wielding a sword, to punish the nations.

On a more prosaic note, in pre-industrial times, the horse signified mobility (it was the quickest way to get around), abundance (it was costly to maintain a horse and develop skills in horsemanship), and power (the horse was effective in battle). Ramon Lull called the horse the ``most noble of beasts''.

Ramon Lull is an interesting figure in the history of Chivalry. He was a knight born to wealth, who lived a rather dissolute life. Later in life, he had a ``conversion'' and he devoted his life to philosophy and God. His biographer claimed that Lull was a Hermeticist, a judgment that is certainly compatible with his written works. In \emph{Perspectives on Initiation}, Guenon tells us:

\begin{quotex}
Within a single organization, a kind of double hierarchy can exist, especially when the apparent leaders are themselves unaware of any link to a spiritual center. In such cases there may exist beside the visible hierarchy made up by those apparent leaders, an invisible hierarchy of which the members may not fulfil any `official' function but who, by their presence alone, nonetheless assure an effective liaison with this center. In the more exterior organizations these representatives of the spiritual centers obviously need not reveal themselves as such ... \end{quotex}


Interestingly enough, these ``invisible'' spiritual leaders often worked as jugglers and \textbf{horse traders}. One can certainly speculate that Lull came into contact with one of these so-called horse traders, from whom he learned Hermetic science.

Finally, this is what Lull wrote about the Knight's horse in his \emph{Book of Knighthood and Chivalry}, an instruction manual for young knights written in his later years:

\paragraph{The Knight's Horse: Noblesse of Courage}

\begin{quotex}
To a knight is given a horse, and also a courser to signify noblesse of courage. And because he is well horsed and high is why he may be seen to be free from fear. And that is the significance, that he ought to be made ready to do all that which behooves the order of chivalry more than another man would. To a horse is given a bridle. And the reigns of the bridle are given to the hands of a knight because the knight may at his own will hold the course the horse and refrain him. And this signifies that the knight ought to refrain his tongue, and hold that he speaks neither foul nor false. And it signifies that he ought to refrain his hands, that he gives not so much that he is suffrous and needy. And that he begs or asks not; nor ought he be so hardy that in his hardiness he has reason and temperance.
\end{quotex}



\flrightit{Posted on 2010-03-06 by Cologero}

\begin{center}* * *\end{center}

\begin{footnotesize}
\begin{sffamily}
\texttt{Vijay Kumar on 2010-03-06 at 22:51 said:}

When do we expect coming of Bhagavan Kalki... an \textit{Avatar}\footnote{\url{http://www.godrealized.org/Avatar.html}} of present era... a messiah competent to uplift Dharma (righteousness)! In times of strife... when mutual trust amongst each other reaches its lowest ebb... coming of a messiah (Avatar) gets necessitated! Even man gods like Mahavira, Gautama Buddha, Jesus Christ or Prophet Mohammed fail in such circumstances!

\textit{Who shall be Bhagavan Kalki}\footnote{\url{http://www.vijaykumar.org/who_shall_be_bhagavan_kalki.html}} is the biggest question! Coming of a man god coupled with powers of Chanakya (most able administrator in history of mankind) is only solution! Humanity looks forward to such a one... an ordinary mortal that against all odds upholds Dharma (righteousness) and succeeds in re-establishing peace tranquility abound!

\hfill

\texttt{GFJF on 2010-03-09 at 05:12 said:}

My thanks. An interesting fact I ran across in Dumezil's `Archaic Roman Religion' is that it was forbidden alike for the brahman and the roman priest of Jupiter to touch horses. According to Dumezil the horse was the animal of Mars and Indra (what relation to Agni?).

\hfill

\texttt{GFJF on 2010-03-09 at 05:36 said:}

Ah, I see they are closely related.

\url{http://www.sacred-texts.com/hin/rigveda/rv06060.htm}


\hfill

\texttt{Max on 2015-02-14 at 07:38 said:}

That priests are not to touch horses are significant. There are two possible interpretations as to the symbolism of the roman legionnaires helmet crests: the mane of a boar or the mane of a horse. In Norse mythology Freyr's boar is called Gullinbursti (Golden mane):

``... to Freyr he gave the boar, saying that it could run through air and water better than any horse, and it could never become so dark with night or gloom of the Murky Regions that there should not be sufficient light where he went, such was the glow from its mane and bristles.''

The priestly equivalent of a knights horse is the boar. Guenon in his essay ``The Wild Boar and the Bear'' relates the symbol of the boar directly to the primordial tradition explaining how Borea means ``land of the wild boar''. If a previous boar-symbolism is replaced by horse-symbolism it may mark a shift of focus from sacerdotal to martial function. The original meaning of wearing the helmet crests could then be interpreted as ``Guardians of the Holy Land''. In ancient times horses ran free as the wind across Eurasia, untouched by men, which means desire was pure.

\end{sffamily}
\end{footnotesize}

